\documentclass[10pt,a4paper]{amsart}
\usepackage[margin=19mm]{geometry}
\usepackage{amsmath,amssymb,mathtools}
\usepackage[scr=rsfs,cal=euler]{mathalfa}
\usepackage{xfrac}
\usepackage[unicode,hidelinks,bookmarks=false]{hyperref}
\newcommand{\ie}{i.e.\ }
\newcommand{\cf}{cf.\ }
\newcommand{\resp}{respectively\ }
\newcommand{\Subsection}{\subsection}
\providecommand{\nobreakdash}{\nobreak\hbox{-}}
\setlength{\emergencystretch}{2em}
\multlinegap0pt
\usepackage{mathtools}
\usepackage{amssymb}
\usepackage{tikz}
\usepackage[shortlabels]{enumitem}
\usepackage{csquotes}
\usepackage{bm}
\newtheorem{theorem}{Theorem}
\newtheorem{lemma}{Lemma}
\def\E{\mathbb E}
\def\R{\mathbb R}
\def\leq {\leqslant}
\title{Helson's conjecture for smooth numbers}
\author{Seth Hardy, Max Wenqiang Xu}
\date{Source: arXiv:2511.03430v3 (CC BY 4.0)}
\begin{document}
\maketitle

\noindent\emph{Source and licence.} Helson's conjecture for smooth numbers. Version arXiv:2511.03430v3 (CC BY 4.0), distributed by arXiv under the Creative Commons Attribution 4.0 International licence, http://creativecommons.org/licenses/by/4.0/, as recorded in the arXiv licence metadata for this exact version and shown on the abstract page https://arxiv.org/abs/2511.03430v3. Full licence notice: \texttt{licenses/arxiv-2511.03430-CC-BY-4.0.txt}.\par\smallskip

\noindent\textit{Editorial scope.} Counted unit: the article's theorem quantitative (source0.tex lines 189--200). The article's complete original proof of it is delivered with it (source0.tex lines 534--668, one \texttt{proof} environment of 135 lines, which the article places away from the statement). No mathematics is written, repaired or re-derived: the statement and the proof are the article's own lines, in the article's own notation, and no retained line is substituted or re-flowed.  Retained uncounted as the source-local context the proof needs: lines 223--225; lines 372--381; lines 386--391; lines 396--402; lines 437--443.  Nothing was omitted from any retained span, so the package carries no omission record.  The retained lines cite 1 works, whose entries are transcribed verbatim from the article's own bibliography source in the e-print and delivered with short editorial labels recorded in the item's reference mapping.  No retained line contains a figure environment, an image-inclusion command or drawing code, so no image is delivered and the package declares no asset dependency.  Editorial formatting only, in addition: an AMS \texttt{amsart} preamble that restates the article's own declarations and macros used by the retained lines, namely the environments \texttt{theorem}, \texttt{lemma} on separate counters, together with the standard packages those lines need.  The retained environments print in this document as Theorem 1, Lemma 1 in order of appearance, whereas the article prints the same items with the numbering of its own full document.  The complete native source of the article as distributed in the arXiv e-print for this exact version is bundled unchanged as \texttt{sources/188501/source0.tex}.
\begin{equation}\label{equ:rankin to true count comparison}
x^{\alpha} \zeta(\alpha,y) \ll \Psi(x,y) \sqrt{\log x \log y},
\end{equation}

\begin{lemma}[Saddle point estimate 1]\label{l:saddle point approx}  
For $2 \log x \leq y \leq x$, say, $u= \frac{\log x}{\log y}$ and for $\alpha(x, y)$ the saddle point, we have
\[
\alpha (x,y) = 1 - \frac{\log (u \log (u + 1))}{\log y} + O \biggl( \frac{1}{\log y} \biggr).
\] 
In fact, for $2 \log x \leq y \leq x^{1/10}$, say, we have 
\[
\alpha (x,y) = 1 - \frac{\log (u \log u)}{\log y} + O \biggl( \frac{\log \log u}{(\log y) (\log u)} + \frac{\log x}{y \log y}\biggr).
\] 
\end{lemma}

\begin{lemma}[Saddle point estimate 2]\label{l:saddle point approx uniform}  
Uniformly for $2 \leq y \leq x$, we have
\[
\alpha (x,y) = \frac{\log ( 1 + y/\log x )}{\log y} \biggl( 1 + O \biggl( \frac{\log \log (1+y)}{\log y} \biggr) \biggr).
\] 
\end{lemma}

\begin{lemma}[Explicit smooth count]\label{l:smooth count in terms of saddle point}
Uniformly for $x \geq y \geq 2$, the number of $y$ smooth numbers up to $x$ satisfies
\[
\Psi(x, y)=\frac{x^\alpha \zeta(\alpha, y)}{\alpha \sqrt{2 \pi(1+(\log x) / y) \log x \log y}}\left(1+O\left(\frac{1}{\log (u+1)}+\frac{1}{\log y}\right)\right),
\]
where $\alpha = \alpha (x,y)$ is the saddle point.
\end{lemma}

\begin{lemma}[Main expectation result]\label{l:exp result moderate y}
Let $f$ be a Steinhaus random multiplicative function, let $\beta$ be a real number such that $|\beta| \leq 2$, say. Suppose also that $t \in \R$, and $0 < \sigma \leq 2$, then we have
\[
\E |F_y (\sigma / 2 + it)|^{\beta} = |\zeta (\sigma, y)|^{\beta^2/4} \exp \biggl( O \biggl( \sum_{p \leq y} \frac{1}{p^{3 \sigma / 2} | 1 - p^{-\sigma/2}|} \biggr) \biggr).
\]
where the implied constant is absolute.
\end{lemma}

\begin{theorem}[Moderate $y$]\label{Thm: main quantitative}
Let $u = \frac{\log x}{\log y}$ and fix $\varepsilon > 0$. For $(\log x)^{1 + \varepsilon} \leq y \leq x^{\frac{1}{(\log \log x)^{1 + \varepsilon}}}$, we have
\begin{equation}\label{equ: alt statement}
\E \Bigl| \sum_{\substack{n \leq x \\ P(n) \leq y}} f(n) \Bigr| \ll \sqrt{\Psi(x,y) \exp \bigl( -u (\log 2 + o(1) \bigr)},
\end{equation}
where the $o(1)$ term goes to zero uniformly for $y$ in this range. This result is derived from the following bound: In the larger range $1000 \log x \leq y \leq x^{\frac{1}{\log \log x}} $, we have
\[
\E \Bigl| \sum_{\substack{n \leq x \\ P(n) \leq y}} f(n) \Bigr| \ll \Psi(x^2, y)^{1/4} \exp \biggl( O \biggl( u \Bigl( \frac{\log x}{y} \Bigr)^{1/3} \biggr) \biggr) (\log x)^{9/8} (\log y)^{1/8},
\]
where all implied constants are uniform for $y$ in this range. As is shown in the proof (see~\eqref{equ: comp to trivial bound}), this
gives better than squareroot cancellation (saving a factor of at least $e^{cu}$ over the trivial bound $\sqrt{\Psi(x,y)}$, for some fixed $c>0$) uniformly on the entire range $C \log x \leq y \leq x^{\frac{1}{3 \log \log x}} $ for some large constant $C \geq 1000$.
\end{theorem}

\begin{proof}[Proof of Theorem~\ref{Thm: main quantitative}]

To begin, suppose that $1000 \log x \leq y \leq x$. By Perron's formula~\cite[Corollary~5.3]{MV2007}, we have
\[
\sum_{\substack{n \leq x \\ P(n) \leq y}} f(n) = \frac{1}{2 \pi i} \int_{1 - i x}^{1 + ix} F_y (s) x^s \frac{ds}{s} + O \Biggl( 1 + \sum_{\substack{x/2 < n \leq x \\ n \neq x}} \min \biggl\{ 1, \frac{1}{|x-n|} \biggr\} + \prod_{p \leq y} \biggl( 1 - \frac{1}{p} \biggr)^{-1} \Biggr).
\]
It follows from straightforward calculations (using Mertens' Theorem to bound the Euler product) that the error term here is $\ll \log x$. 

Since the integrand is analytic and has no poles apart from $s=0$, we can shift the integral to the line $\Re(s) = \alpha(x^2, y)/2$. Specifically, we have
\begin{multline*}
\frac{1}{2 \pi i} \int_{1 - i x}^{1 + ix} F_y (s) x^s \frac{ds}{s} = \frac{1}{2 \pi i} \int_{\alpha (x^2, y) - ix}^{\alpha (x^2, y) + ix} F_y (s) x^s \frac{ds}{s} \\  
+ O \biggl( \frac{1}{x} \int_{\alpha (x^2, y)/2}^1 \bigl( |F_y (\sigma / 2 + ix)| + |F_y (\sigma / 2 - ix)| \bigr) x^{\sigma} \, d \sigma \biggr) .
\end{multline*}
For the remainder of this section, we let $\alpha \coloneqq \alpha (x^2, y)$. Not carefully that, in the following sections, we will instead take $\alpha$ to be the typical saddle point corresponding to $x$ and $y$. By the triangle inequality, it follows from the previous equations that 
\begin{equation}\label{eqn: whole}
   \begin{split}
       \E \Bigl| \sum_{\substack{n \leq x \\ P(n) \leq y}} f(n) \Bigr| \ll  x^{\alpha(x^2, y)/2} \int_{-x}^{x} \E \biggl| \frac{F_y (\alpha / 2 + it)}{\alpha / 2 + it} \biggr| \, dt + \log x \\ 
+ \frac{1}{x} \int_{\alpha/2}^1 \bigl( \E |F_y (\sigma / 2 + ix)| + \E |F_y (\sigma / 2 - ix)| \bigr) x^{\sigma} \, d \sigma. 
   \end{split} 
\end{equation}
We now apply Lemma~\ref{l:exp result moderate y} to each expectation here. For the last term, we note that the bound is maximised on our range of integration when $\sigma = \alpha/2$, so that we have $\E |F_y (\sigma / 2 + it)| \ll \zeta(\alpha, y)^{1/4} \exp \bigl( O \bigl( \sum_{p \leq y} \frac{1}{p^{3 \alpha / 2} | 1 - p^{-\alpha /2}|} \bigr) \bigr)$. Then using the fact that $\int_{-x}^{x} \frac{dt}{|\alpha/2 + it|} \ll \log x$ uniformly on our range of $\alpha = \alpha (x^2, y)$ (see Lemma~\ref{l:saddle point approx}), we find that the first term dominates the last term, and putting this all together, we deduce that
\begin{equation}\label{equ: moderate y before triangle}
\E \Bigl| \sum_{\substack{n \leq x \\ P(n) \leq y}} f(n) \Bigr| \ll x^{\alpha (x^2, y)/2} \zeta \bigl( \alpha (x^2, y), y \bigr)^{1/4} \exp \biggl( O \biggl( \sum_{p \leq y} \frac{1}{p^{3 \alpha / 2} |1 - p^{-\alpha / 2}|} \biggr) \biggr) (\log x).
\end{equation}
Here we have also used the fact that our bound certainly exceeds $\log x$, so the first term dominates the right-hand side of \eqref{eqn: whole}.

We proceed to bound the prime number sum in the exponent. For $\alpha \log p \leq 1$, we have the lower bound $|1 - p^{-\alpha / 2}| \gg \alpha \log p$, and for $\alpha \log p \geq 1$, we have $|1 - p^{-\alpha / 2}| \gg 1$. When $x$ is sufficiently large, we certainly have $\alpha (x^2, y) \log y \geq 1$ for $y \geq 1000 \log x$ by Lemma~\ref{l:saddle point approx uniform}, so we deduce that
\begin{equation}\label{eqn: primesum}
\sum_{p \leq y} \frac{1}{p^{3 \alpha / 2} |1 - p^{-\alpha / 2}|} \ll \sum_{p \leq e^{1/\alpha}} \frac{1}{\alpha p^{3 \alpha / 2} \log p} + \sum_{e^{1 / \alpha} < p \leq y} \frac{1}{p^{3 \alpha / 2}} . 
\end{equation}
Suppose now that $y \geq (\log x)^{4}$, then we certainly have $\alpha (x^2, y) \geq 7/10$, say, hence $3 \alpha (x^2, y) / 2 \geq 1.05$. It then follows from the above that
\[
\sum_{p \leq y} \frac{1}{p^{3 \alpha / 2} |1 - p^{-\alpha / 2}|} \ll 1
\]
Now suppose that $1000 \log x \leq y \leq (\log x)^4$. Then we know by Lemma~\ref{l:saddle point approx uniform} that $\frac{1}{\log y} \leq \alpha (x^2, y) \leq 4/5$ when $x$ is sufficiently large, and it follows from dyadic decomposition that
\[
\sum_{p \leq e^{1/\alpha}} \frac{1}{\alpha p^{3 \alpha / 2} \log p} \ll 1 + \frac{1}{\alpha} \sum_{1 \leq i \leq 1/\alpha} \frac{e^{i(1 - 3 \alpha / 2)}}{i^2} \ll e^{1/\alpha}  ,
\]
say. In this range of $y$, i.e., $1000 \log x \leq y \leq (\log x)^4$, we also have
\[
\sum_{e^{1 / \alpha} < p \leq y} \frac{1}{p^{3 \alpha / 2}} \ll 1 + \frac{y^{1 - 3 \alpha / 2}}{\log y} \]
 Hence in view of \eqref{eqn: primesum} and combining the above estimates, we deduce that
\[
\sum_{p \leq y} \frac{1}{p^{3 \alpha / 2} |1 - p^{-\alpha / 2}|} \ll e^{1/\alpha} + \frac{y^{1 - 3 \alpha / 2}}{\log y} ,
\]
uniformly for all $1000 \log x \leq y \leq x$. Inserting this estimate into~\eqref{equ: moderate y before triangle}, we have
\[
\E \Bigl| \sum_{\substack{n \leq x \\ P(n) \leq y}} f(n) \Bigr| \ll x^{\alpha (x^2, y) / 2} \zeta \bigl( \alpha (x^2, y), y \bigr)^{1/4} \exp \biggl( O \biggl( e^{1/\alpha} + \frac{y^{1 - 3 \alpha / 2}}{\log y} \biggr) \biggr) (\log x) .
\]
Now applying equation~\eqref{equ:rankin to true count comparison} (which follows from Lemmas~\ref{l:smooth count in terms of saddle point} and~\ref{l:saddle point approx}), we have
\begin{equation}\label{equ: quant bound in terms of Psi (x^2, y)}
\E \Bigl| \sum_{\substack{n \leq x \\ P(n) \leq y}} f(n) \Bigr| \ll \Psi(x^2, y)^{1/4} \exp \biggl( O \biggl( e^{1/\alpha} + \frac{y^{1 - 3 \alpha / 2}}{\log y} \biggr) \biggr) (\log x)^{9/8} (\log y)^{1/8}.
\end{equation}
We now bound the error terms in the exponential more explicitly. In particular, we will show that
\[
e^{1/\alpha} + \frac{y^{1 - 3 \alpha / 2}}{\log y} \ll 1 + u \biggl( \frac{\log x}{y} \biggr)^{1/3}. 
\]
For this, we will make use of Lemma~\ref{l:saddle point approx uniform}, which (for ease of reference) tells us that
\begin{equation}\label{equ: saddle point uniform}
\alpha (x,y) = \frac{\log (1 + y/\log x)}{\log y} \biggl( 1 + O \biggl( \frac{\log \log (1+y)}{\log y} \biggr) \biggr)
\end{equation}
uniformly for $2 \leq y \leq x$. Note that we wish to apply this with $\alpha = \alpha(x^2,y)$, and in particular it implies that $\alpha (x^2, y) \geq \frac{8 \log (1 + y / 2 \log x)}{9 \log y}$, say, so for $1000 \log x \leq y \leq x$, we have
\[
y^{-3\alpha / 2} \ll \exp \biggl( - \frac{4}{3} \log \biggl(1 + \frac{y}{2\log x}\biggr) \biggr) \ll \biggl( \frac{\log x}{y} \biggr)^{4/3}.
\]
From this it follows that
\[
\frac{y^{1 - 3 \alpha / 2}}{\log y} \ll u \biggl( \frac{\log x}{y} \biggr)^{1/3}.
\]
Furthermore, by~\eqref{equ: saddle point uniform}, we have
\[
e^{1/\alpha} \ll \exp \biggl( \frac{2 \log y}{\log (1 + y/2\log x)} \biggr).
\]
In the range $y \geq (\log x)^2$, say, this is certainly $\ll 1$, and for $1000 \log x \leq y \leq (\log x)^2$, this satisfies the bound $\ll y^{3/10}$. These estimates together leads to that the bound
\[
e^{1/\alpha} \ll 1 +  u \biggl( \frac{\log x}{y} \biggr)^{1/3},
\]
holds uniformly for $1000 \log x \leq y \leq x$. Now inserting these estimates into~\eqref{equ: quant bound in terms of Psi (x^2, y)}, we have 
\begin{equation}\label{equ: L1 norm in terms of psi(x^2, y) bound}
\E \Bigl| \sum_{\substack{n \leq x \\ P(n) \leq y}} f(n) \Bigr| \ll \Psi(x^2, y)^{1/4} \exp \biggl( O \biggl( u \Bigl( \frac{\log x}{y} \Bigr)^{1/3} \biggr) \biggr) (\log x)^{9/8} (\log y)^{1/8},
\end{equation}
This proves the second formulation in Theorem~\ref{Thm: main quantitative} on the range $1000 \log x \leq y \leq x$, though as we will see, the result is only non-trivial in a more restrictive range of $y$, so it is not stated for this range in the theorem. To elucidate our saving, we now undo the application of Lemma~\ref{l:smooth count in terms of saddle point} so that we can compare the right-hand side to the trivial bound, $\sqrt{\Psi(x,y)}$. In particular, we will show that in the range $1000 \log x \leq y \leq x^{\frac{1}{\log \log x}}$, we have
\[
\Psi (x^2, y) \ll \frac{\Psi(x,y)^2}{\log y} \exp \biggl( - u \biggl( 2 \log 2 + O \biggl( \Bigl( \frac{\log x}{y} \Bigr)^{1/3} + \frac{\log \log u}{\log u} \biggr) \biggr) \biggr),
\]
and from this we will deduce~\eqref{equ: alt statement}. First of all, by Lemma~\ref{l:smooth count in terms of saddle point}, we have
\begin{equation}\label{equ: order of mag comparrison}
\Psi (x^2, y) \asymp \Psi (x,y)^2 x^{2 (\alpha (x^2, y) - \alpha (x,y))} \frac{\zeta \bigl( \alpha (x^2, y), y \bigr)}{\zeta \bigl( \alpha (x,y), y \bigr)^2} \frac{\alpha(x,y)^2}{\alpha(x^2, y)} \sqrt{\log x \log y}.
\end{equation}
We begin by considering the $x^{2 (\alpha (x^2, y) - \alpha (x,y))}$ term. Applying the second part of Lemma~\ref{l:saddle point approx}, we have
\begin{equation}\label{equ:alpha(x^2,y) vs alpha(x,y)}
\begin{split}
\alpha (x^2, y) - \alpha (x,y) &= \frac{\log (\frac{u \log u}{2 u \log 2u})}{\log y} + O \biggl( \frac{\log \log u}{(\log y)(\log u)} \biggr) \\
& = - \frac{\log 2}{\log y} + O \biggl( \frac{\log \log u}{(\log y)(\log u)} \biggr) .
\end{split}
\end{equation}
on the range $1000 \log x \leq y \leq x^{\frac{100}{\log \log x}}$, say. This is precisely where our saving of the form $\exp (- u \log 2 / 2)$ comes from. We now move on to bounding the quotient of zeta functions. First of all, we note that for any $\sigma > 0$, we have
\[
\log \zeta (\sigma, y) = \sum_{p \leq y} \frac{1}{p^{\sigma}} + O \biggl( \sum_{p \leq y} \frac{1}{p^{2 \sigma} |1 - p^{-\sigma}|} \biggr).
\]
This follows from the fact that $\log (1 - x) = - x + O \bigl( \frac{|x|^2}{1 - |x|} \bigr) $ uniformly for $|x| < 1$. For both $\sigma = \alpha (x,y)$ and $\sigma = \alpha (x^2, y)$, the handling of the error term is identical to the similar term that appeared previously, specifically the error term in the exponential in equation~\eqref{equ: moderate y before triangle}. By following identically the handling of that term (or by noting that the terms in our case are slightly smaller), we can deduce the same bound here, so that 
\[
\frac{\zeta \bigl( \alpha (x^2, y), y \bigr)}{\zeta \bigl( \alpha (x,y), y \bigr)^2} = \exp \biggl( \sum_{p \leq y} \frac{1}{p^{\alpha(x^2, y)}} - 2 \sum_{p \leq y} \frac{1}{p^{\alpha (x,y)}} + O \biggl(1 +  u \Bigl( \frac{\log x}{y} \Bigr)^{1/3} \biggr) \biggr).
\]
We now bound the prime number sums, which we evaluate using~\cite[Lemma~7.4]{MV2007} to find that
\[
\sum_{p \leq y} p^{-\alpha (x, y)} = \frac{y^{1 - \alpha(x,y)}}{(1 - \alpha(x,y)) \log y} +  \log \Bigl( \frac{1}{1 - \alpha(x,y)} \Bigr) + O \biggl( \frac{y^{1 - \alpha(x,y)}}{(1 - \alpha(x,y))^2 (\log y)^2} \biggr).
\]
Applying Lemma~\ref{l:saddle point approx}, we obtain
\[
\sum_{p \leq y} p^{-\alpha (x, y)} = u + \log \Bigl( \frac{1}{1 - \alpha(x,y)} \Bigr) + O \biggl( \frac{u \log \log u}{\log u} + \frac{u \log x}{y}\biggr),
\]
uniformly for $2 \log x \leq y \leq x^{\frac{100}{\log \log x}}$, say. Similarly, in the same range, we have 
\[
\sum_{p \leq y} p^{-\alpha (x^2 , y)} = 2u + 2 \log \Bigl( \frac{1}{1 - \alpha(x^2 ,y)} \Bigr) + O \biggl( \frac{u \log \log u}{\log u} + \frac{u \log x}{y}\biggr).
\]
Therefore, we obtain the bound
\[
\frac{\zeta \bigl( \alpha (x^2, y), y \bigr)}{\zeta \bigl( \alpha (x,y), y \bigr)^2} = \Biggl( \frac{\bigl( 1 - \alpha (x,y) \bigr)^2}{1 - \alpha (x^2 ,y)} \Biggr)  \exp \biggl( O \biggl( u \Bigl( \frac{\log x}{y} \Bigr)^{1/3} + \frac{u \log \log u}{\log u} \biggr) \biggr).
\]
Combining this result with~\eqref{equ: order of mag comparrison} and~\eqref{equ:alpha(x^2,y) vs alpha(x,y)}, we deduce that
\begin{multline*}
\Psi (x^2, y) = \Psi(x,y)^2 \exp \biggl( - u \biggl( 2 \log 2 + O \biggl( \Bigl( \frac{\log x}{y} \Bigr)^{1/3} + \frac{\log \log u}{\log u} \biggr) \biggr) \biggr) \\ \times \biggl( \frac{\alpha(x,y)^2}{\alpha(x^2, y)} \biggr) \Biggl( \frac{\bigl( 1 - \alpha (x,y) \bigr)^2}{1 - \alpha (x^2 ,y)} \Biggr) \sqrt{\log x \log y}.
\end{multline*}
Finally, we note that on the range $2 \log x \leq y \leq x^{\frac{100}{\log \log x}}$, say, Lemma~\ref{l:saddle point approx} implies that $\frac{\alpha(x,y)^2}{\alpha(x^2, y)} \ll \frac{1}{\log y}$, and by Lemma~\ref{l:saddle point approx} we obtain the bound $\bigl( \frac{( 1 - \alpha (x,y) )^2}{1 - \alpha (x^2 ,y)} \bigr) \sqrt{\log x \log y} \ll \sqrt{u} \log u$, which can be engulfed into the error in the exponential. Therefore, as desired, we have
\[
\Psi (x^2, y) \ll \frac{\Psi(x,y)^2}{\log y} \exp \biggl( - u \biggl( 2 \log 2 + O \biggl( \Bigl( \frac{\log x}{y} \Bigr)^{1/3} + \frac{\log \log u}{\log u} \biggr) \biggr) \biggr),
\]
which holds uniformly for $1000 \log x \leq y \leq x^{\frac{1}{\log \log x}}$, say. Inserting this estimate into equation~\eqref{equ: L1 norm in terms of psi(x^2, y) bound} and using the fact that $(\log x)^{9/4} (\log y)^{-1/4} \ll u \log x$, we have
\begin{equation}\label{equ: comp to trivial bound}
\E \Bigl| \sum_{\substack{n \leq x \\ P(n) \leq y}} f(n) \Bigr| \ll \\ 
\sqrt{\Psi(x,y)\exp \Bigl(-u \biggl( \log 2 -  O \biggl(\Bigl( \frac{\log x}{y} \Bigr)^{1/3} + \frac{\log \log u}{\log u} \biggr) \biggr) \biggr) (\log x)^2}.
\end{equation}
uniformly for $1000 \log x \leq y \leq x^{\frac{1}{\log \log x}}$. Now, in the range $(\log x)^{1 + \varepsilon} \leq y \leq x^{\frac{1}{(\log \log x)^{1 + \varepsilon}}}$, say, we have $\frac{\log x}{y} = o(1)$ and $\log \log x = o(u)$, where the $o(1)$ terms goes to zero uniformly for fixed $\varepsilon > 0$. Therefore, equation~\eqref{equ: alt statement} holds on this range of $y$, as required. It is also clear from this inequality that we have better than squareroot cancellation (for smaller $y$) whenever $y \geq C \log x$ where $C \geq 1000$ is a large constant, and for $y$ closer to $x$ we have better than squareroot cancellation whenever $u > \frac{2}{\log 2} \log \log x \approx 2.885 \log \log x$ to overcome the $(\log x)^{2}$ term. Therefore, for some large constant $C \geq 1000$, the range $C \log x \leq y \leq x^{\frac{1}{3\log \log x}}$ is such that we always have a saving of at least $e^{cu}$ for some fixed $c > 0$ on this range. This completes the proof of Theorem~\ref{Thm: main quantitative}.
\end{proof}

\begin{thebibliography}{99}
\bibitem[MV2007]{MV2007} H.~L. Montgomery and R.~C. Vaughan. \newblock {\em Multiplicative number theory. {I}. {C}lassical theory}, volume~97 of {\em Cambridge Studies in Advanced Mathematics}. \newblock Cambridge University Press, 2007.
\end{thebibliography}
\end{document}
